\documentclass[12pt, ukrainian]{article}
\usepackage[a4paper]{geometry}
\setlength\parindent{0mm}
\geometry{top=1.1cm,bottom=1.1cm,left=1.1cm,right=1.1cm}
\pagestyle{empty}
\usepackage[T2A]{fontenc}
\usepackage[utf8]{inputenc}
\usepackage[ukrainian]{babel}
\usepackage{graphicx}
\usepackage{caption}
\usepackage{verbatim, listings, clrscode3e, fancyvrb}
\DeclareCaptionFormat{nocaption}{}
\captionsetup{format=nocaption,aboveskip=0pt,belowskip=0pt}
\usepackage[Export]{adjustbox}
\adjustboxset{max size={0.9\linewidth}{0.9\paperheight}}
\def\code#1{\Verb!#1!}
\begin{document}
{{16.05.2024 Контрольна робота, варіант  \No \textbf { 17 }} група К\_\_ ПІБ {\parbox {2.5 cm} \dotfill}\par}
{
%\begin {large}
\begin{minipage}[t]{9.7cm}
У завданнях 1--3 вказати, що буде виведено за виконання. Повідомлення інтерпретатора про помилку позначати ERROR


               
{\begin{center}Завдання {1}\end{center}}
\vspace{-4mm}

\begin{verbatim}
g = ('a','b','c','d','e',0,1,2,3)
print(g[5])

\end{verbatim}            

\vspace{15mm}
               
{\begin{center}Завдання {2}\end{center}}
\vspace{-4mm}

\begin{verbatim}
f=[10,11,12,13]
del f[0:8]
print(f)

\end{verbatim}            


\end{minipage}
\vline \hspace{8mm}
\begin{minipage}[t]{8.0cm}                
               
{\begin{center}Завдання {3}\end{center}}
\vspace{-4mm}

\begin{verbatim}
c = 0
j = 1

class B:
    j = 2
    
    def __init__(self):
        global j
        self.c = 3
        j = 4
        B.j = 5

obj = B()
try:
    print(c, end=' ')
    print(j, end=' ')
    print(B.j, end=' ')
    print(obj.c)
except:
    print('ERR')

\end{verbatim}            

\end{minipage}


               
{\begin{center}Завдання {4}\end{center}}
\vspace{-4mm}

 Задано числову послідовність \kw{s} довжини 6000. Записати ВИРАЗ, значенням якого є список, що складається з елементів послідовності, що діляться на 13, записаних в обернено\-му порядку.\par

\vspace{10mm}
               
{\begin{center}Завдання {5}\end{center}}
\vspace{-4mm}
Змінні \kw{next} та \kw{prev} вузла списку позначають наступний та попередній за ним вузол відповідно. Починаючи з голови, у список записано послідовні цілі числа від~$-28$ до~$18$. Нехай змінна \kw{p} позначає вузол, що зберігає значення~$-1$.
Яке число зберігається у вузлі \kw{p.prev.prev.next.next.prev} ?\par

               
{\begin{center}Завдання {6}\end{center}}
\vspace{-4mm}
Значенням змінної \kw{a} є цілочисловий масив типу $numpy.ndarray$ розмірів $7{\times}3$. На скільки збільшиться сума елементів масиву \kw{a} після виконання
\begin{verbatim}
a.T[-3] += 3
\end{verbatim}


Якщо код некоректний, то в якості відповіді зазначити ERROR.\par

               
{\begin{center}Завдання {10}\end{center}}
\vspace{-4mm}

Написати функцію $f$, яка для файлу, заданого шляхом, розв’язує задачу:\\ Текстовий файл містить послідовність чисел $a_1$, $a_2$, $\ldots$, $a_t$, записану у форматі одне число на рядок. Знайти третє найбільше значення, що записане у файлі. Кожне значення, що записане у файлі, враховується тільки один раз. Для послідовності -1, -4, -12, -5, -1, 0 таким значенням буде -4.

В усіх випадках розміри файлів такі, що їх вміст повністю завантажений у пам'ять бути не може. Якість коду також оцінюється.



\vspace{10mm}\par

\newpage
\begin{minipage}[t]{9.7cm}
               
{\begin{center}Завдання {7}\end{center}}
\vspace{-4mm}

\begin{center}
	\adjustimage{max size={0.8\linewidth}{0.8\paperheight}}
	{../15BinTree/trees_exp/34.png}
\end{center}
Указати послідовність міток вершин дерева, зображеного на рис., що утвориться в результаті обходу дерева в оберненому порядку. Мітки записувати через пробіл.\par Алгоритм починає роботу з кореня (на рис. це верхня вершина).\par

\end{minipage}
\vline \hspace{8mm}
\begin{minipage}[t]{8.0cm}
               
{\begin{center}Завдання {8}\end{center}}
\vspace{-4mm}
Для графа на рисунку з вершини 3 запущено алгоритм пошуку в ширину, що будує кістякове дерево. Списки суміжності вершин переглядаються алгоритмом за зростанням міток вершин.\par
\begin{center}
	\adjustimage{max size={0.9\linewidth}{0.9\paperheight}}
	{../16Graphs/Traversals/228.png}
\end{center}
\par Які з наведених ребер входять до складу отриманого кістякового дерева?\par
1) (1, 3)\par
2) (0, 1)\par
3) (1, 4)\par
4) (4, 6)\par
5) (0, 6)\par
6) (0, 2)\par
{\vskip 15pt} \par

\end{minipage}

               
{\begin{center}Завдання {9}\end{center}}
\vspace{-4mm}
Рядки відомості через кому містять порядковий номер (ціле), назву предмету (зна\-ків пунктуації не містить), прізвище студента, ім'я студента, оцінку в балах за 100 бальною шкалою.

Білі символи навколо роздільників не допускаються.

Приклад такого рядка присвоєно змінній \kw{s}.



\begin{verbatim}
s = '13,algebra,surname,name,77'
\end{verbatim}


Нижче наведено код, за виконання якого значенням змінної \kw{out} має стати прізвище студента.

Для наведеного прикладу значенням змінної \kw{out} має стати рядок \code{surname}



\begin{verbatim}
res = re.fullmatch(r'XXX', s)
s = ''
out = YYY
\end{verbatim}


Код має працювати не тільки для конкретного рядка з прикладу, але й для кожного рядка з відомості.


Який рядок треба записати на місце \kw{XXX} ?
Який рядок треба записати на місце \kw{YYY} ?\par

%\end {large}
}
%end
\end{document}

